\section{Löwenheim–Skolem 与非标准模型}


\subsection{紧致性制造大模型}


在 PA 的语言中加入一组常元 $\{c_s:s\in S\}$，并加入

\[
c_s\ne c_t\qquad(s\ne t).
\]

任意有限子集只要求有限多个常元两两不同，可以在标准自然数中解释；由紧致性，整个扩充理论有模型。因此 PA 有至少 $|S|$ 个元素的模型。取不可数的 $S$，就得到不可数的 PA 模型。

这说明一阶 PA 不能在同构意义下唯一刻画标准自然数；非标准模型中会出现标准部分以外的元素。

\subsection{向下 Löwenheim–Skolem}


若一个一阶理论使用可数语言，并且有无限模型，则向下 Löwenheim–Skolem 定理保证它有可数模型。对实闭域或集合论的相关一阶理论，也可得到可数模型；这里说的是理论在某个可数结构中成立，不是说外部的实数集合本身变成可数。

\subsection{Skolem 悖论的层次区别}


一个可数模型可以满足“没有从模型内部枚举出全部对象的函数”这类内部不可数性句子。模型内部的“可数”是指存在模型中的某个函数或集合；外部观察者看到的模型底层集合是否可数，是另一种元理论性质。两者不矛盾。

因此“可数模型满足不可数性”不是逻辑错误，而是内部量词与外部集合大小混用造成的表面悖论。
